% ID: SEL-D0071-OB-E02
% ESTADO: validada. S071.
\textbf{a)} La función es $f(x)=|x-1|$. Sus dos semirrectas forman una V con vértice en $(1,0)$. Pasa también por $(0,1)$ y $(2,1)$.
\begin{center}
\setlength{\unitlength}{0.9mm}
\begin{picture}(97,56)(-29,-9)
\qbezier(-24,0)(19,0)(62,0)
\qbezier(0,-4)(0,19)(0,42)
\qbezier(-20,39)(0,19.5)(20,0)
\qbezier(20,0)(40,19.5)(60,39)
\put(63,0){$x$}\put(2,42){$f(x)$}
\put(15,-6){\scriptsize $(1,0)$}
\put(2,21){\scriptsize $(0,1)$}
\put(42,14){\scriptsize $(2,1)$}
\end{picture}
\end{center}
\textbf{b)} Cada rama lineal es continua en su interior. En $1$, ambos límites laterales y $f(1)$ son cero. Por tanto, es continua en todo $\mathbb R$.

\textbf{c)} Las pendientes laterales en $1$ son $-1$ y $1$. Como no coinciden, la función es derivable en $\mathbb R\setminus\{1\}$, con derivada $-1$ a la izquierda y $1$ a la derecha.
